\chapter{Один train scale и dimensionless blind-scorecard}

RG-аудит доказал неизбежность одного размерного calibration input.
Зафиксируем его без выбора конкретной наблюдаемой:
\begin{equation}
 m_{\mathrm{ref}}
 :=m_f=\chi.
\end{equation}
Это определение устанавливает единицу массы, но не идентифицирует finite
fermion с мюоном, \(\tau\)-лептоном, нейтрино или новым hidden state.

\section{Пререгистрированные dimensionless outputs}

После одной scale-калибровки parent action даёт:
\begin{align}
 \frac{m_{f,1}}{m_{\mathrm{ref}}}
 &=
 \frac{m_{f,2}}{m_{\mathrm{ref}}}=1,
 \\
 \frac{m_{s,a}^2}{m_{\mathrm{ref}}^2}
 &=4,
 \qquad a=1,2,3,
 \\
 \frac{m_A^2}{m_{\mathrm{ref}}^2}
 &=3,
 \\
 g^2(\mu_{\mathrm{spec}})
 &=\frac38,
 \\
 b_{U(1)}
 &=2,
 \\
 |\sin\phi_{\mathrm{vac}}|
 &=1.
\end{align}
Дополнительно остаются две CP-сопряжённые ветви
\(\phi_{\mathrm{vac}}=\pm\pi/2\), пока orientation-odd selector не
построен.

\begin{proposition}[One-scale mathematical prediction set]
После одного train scale модель содержит несколько parameter-free
dimensionless следствий: mass ratios, degeneracies, matching coupling,
one-loop beta coefficient и максимальную CP-фазу. Они являются
математическими outputs parent action.
\end{proposition}

\section{Что не является blind prediction}

Величины
\begin{equation}
 \pi^2+2\pi+\frac23,
 \qquad
 23+\pi^{-1}
\end{equation}
использовались при проектировании role-graded двухсекторного Hessian gate.
Они проверяют внутреннюю совместимость общей метрики, но не могут
считаться независимыми blind observables версии III.

Аналогично равные edge-модули и maximal phase использовались как критерии
выбора flattening potential. Их последующие Hessian masses являются
новыми outputs, но сам факт выбранного vacuum не является внешним blind
test.

\section{Readout gap}

Для физической проверки требуется заранее заданная карта
\begin{equation}
 \begin{aligned}
 \mathcal R_{\mathrm{phys}}:\quad
 &\{\text{finite eigenmodes, collective coordinates}\}
 \\
 &\longrightarrow
 \{\text{наблюдаемые particle states и couplings}\}.
 \end{aligned}
\end{equation}
Текущий parent action фиксирует left/right bimodule charges и spectrum, но
не задаёт:
\begin{enumerate}
 \item соответствие двух Dirac-пар поколениям charged leptons;
 \item соответствие трёх scalar modes Higgs, radion или hidden scalars;
 \item физическую идентичность relative-\(U(1)\) vector;
 \item neutrino mass operator и его связь с collective norm
 \(23+\pi^{-1}\).
\end{enumerate}
Без этой карты сравнение mass ratios с экспериментом является
post-hoc assignment.

\begin{nogocriterion}[Readout-before-data rule]
Ни один finite mode не разрешено называть конкретной известной частицей
после просмотра его предсказанной массы. Representation, charges и
interaction vertices должны однозначно фиксировать идентификацию до
численного сравнения.
\end{nogocriterion}

\section{Контрольная charged-lepton идентификация}

Если условно отождествить две Dirac-пары с двумя различными
charged-lepton поколениями, модель предсказывает точное вырождение
\begin{equation}
 \frac{m_{\ell_2}}{m_{\ell_1}}=1.
\end{equation}
Известные charged-lepton поколения невырождены. Следовательно, такая
прямая readout-карта закрывается без тонкой численной проверки.

\begin{proposition}[Direct charged-lepton readout no-go]
Минимальные две Dirac-пары не могут непосредственно представлять два
различных наблюдаемых charged-lepton поколения. Для их использования
нужны дополнительный family operator, radiative splitting или иная
representation map, отсутствующая в текущей версии.
\end{proposition}

\section{Neutrino blind gate}

Collective norm
\begin{equation}
 N_\nu=23+\pi^{-1}
\end{equation}
не является eigenvalue finite Dirac operator и после аудита тома II не
имеет статуса физического neutrino denominator. Версия III пока не
построила Majorana/Dirac neutrino mass matrix, selector поколений и
readout для mass splittings. Поэтому neutrino sector не предоставляет
второй независимый blind test.

\begin{theorem}[Blind-scorecard verdict]
One-scale mathematical scorecard непуст, но физический двухсекторный
scorecard ещё не определён. Прямая charged-lepton идентификация закрывается
из-за exact degeneracy, а neutrino readout отсутствует. Следовательно,
Definition of Done версии III пока не пройден, несмотря на существенное
замыкание parent action.
\end{theorem}

\section{Текущий scorecard}

\begin{center}
\begin{tabular}{p{0.30\textwidth}p{0.20\textwidth}p{0.30\textwidth}}
\toprule
Пункт & Статус & Причина \\
\midrule
один train scale & pass & \(m_{\mathrm{ref}}=\chi\) \\
dimensionless model outputs & pass & ratios \(1,2,\sqrt3\), \(g^2=3/8\), \(b=2\) \\
charged-lepton blind test & fail & direct map даёт exact degeneracy \\
neutrino blind test & open & mass operator/readout отсутствует \\
двухсекторный Definition of Done & fail & нет двух физических blind observables \\
\bottomrule
\end{tabular}
\end{center}

\begin{protocol}[Следующий гейт]
До новых численных сравнений необходимо вывести representation-level
readout:
\begin{enumerate}
 \item определить, являются ли finite modes hidden states или Standard
 Model states;
 \item если это charged leptons, построить family-splitting operator из
 уже объявленной симметрии;
 \item построить neutrino mass operator, не используя прежний denominator
 как target;
 \item только затем заморозить два blind dimensionless observables.
\end{enumerate}
\end{protocol}

Машинный scorecard сохранён в
\path{s2t_v3_one_scale_blind_scorecard_results.json}.